\hypertarget{link_to_pictlieee}{}{\Large \bf pictlieee}

{\it Proportional-Integral (PI) controller with non-windup limit on the integral term, compliant with IEEE standards}.

\begin{center}
\vspace*{-2ex}
\includegraphics{codegen/pictlieee.pdf}
\vspace*{-2ex}
\end{center}

This PI controller involves a limiter and a non-windup integrator compliant with the IEEE specifications in:
\begin{itemize}
\item IEEE recommended practice for excitation system models for power system stability studies, IEEE Std 421.5-1992
\item IEEE recommended practice for excitation system models for power system stability studies, IEEE Std 421.5-2005 (Revision of IEEE Std 421.5-1992)
\item IEEE recommended practice for excitation system models for power system stability studies, IEEE Std 421.5-2016 (Revision of IEEE Std 421.5-2005)
\end{itemize}

The IEEE standard specfies that the $x_1$ variable is frozen as soon as $x_j$ reaches its (lower or upper) limit. This is done by merely setting the input of the integrator to zero.

However, the standard does not specify the condition under which $x_1$ is let to vary again, i.e. when the integrator is put back into service. A simple way would be to re-activate the integrator as soon $x_j$ gets back within its limits. However, at that time instant, the $x_k$ variable may have a value such that $x_j = K_p x_k + x_1$ is pushed again to its (lower or upper) limit. The system would then be trapped into a limit cycle (which would not match the behaviour of the system to model !).

With the implementation shown in the above block diagram, the limit cycle is avoided by relying on the upper integrator to decide when the lower integrator can be released. The upper integrator is used to observe what the evolution of the system would be with $x_1$ free to vary: the lower integrator is re-activated only when $K_p x_k + x_2$ is in  $[x_j^{min} \; x_j^{max}]$. As it is in ``open loop'', the additional integrator does not interact with the rest of the system.

\vspace{2ex}
Syntax : \hspace*{10ex}\parbox[t]{\textwidth}{\sf 
\&  pictlieee \\
name of variable $x_k$\\
name of variable $x_j$\\
data name, parameter name or math expression for $K_i$\\
data name, parameter name or math expression for $K_p$\\
data name, parameter name or math expression for $x_j^{min}$\\
data name, parameter name or math expression for $x_j^{max}$
}

~ \hrule ~

Internal states: $x_1$ and $x_2$ 

Discrete variables : $z \in \{-2,-1,0,1,2\}$

Equations :
\begin{equation*}
\mbox{if} \;\; z = 0 \mbox{\; : \;}
\begin{cases}
0 & =  K_p x_k + x_1 - x_j \\
\dot{x}_1 & =  K_i \; x_k  \\
0 & =  x_2 - x_1 
\end{cases}
\end{equation*}
\begin{equation*}
\mbox{if} \;\; z = 2 \mbox{\; : \;}
\begin{cases}
0 & =  x_j^{max} - x_j  \\
\dot{x}_1 & =  0  \\
0 & = x_2 - x_1 
\end{cases}
\mbox{\;\;\;\;\;\;\;\;\;\;\;\;\;\;if} \;\; z = 1 \mbox{\; : \;}
\begin{cases}
0 & =  x_j^{max} - x_j \\
\dot{x}_1 & = 0  \\
\dot{x}_2 & = K_i \; x_k  
\end{cases}
\end{equation*}
\begin{equation*}
\mbox{if} \;\; z = -2 \mbox{\; : \;}
\begin{cases}
0 & =  x_j^{min} - x_j  \\
\dot{x}_1 & =  0  \\
0 & = x_2 - x_1 
\end{cases}
\mbox{\;\;\;\;\;\;\;\;\;\;\;\;\;\;if} \;\; z = -1 \mbox{\; : \;}
\begin{cases}
0 & =  x_j^{min} - x_j \\
\dot{x}_1 & = 0  \\
\dot{x}_2 & = K_i \; x_k  
\end{cases}
\end{equation*}

\pagebreak
Discrete transitions (note the use of $x_2$ in the tests marked with (*) ):

\hspace*{20ex}\begin{minipage}[t]{15cm}
\begin{algorithmic} 
\IF {$z = 0$}
  \IF {$x_j > x_j^{max}$} 
    \STATE $z \gets 2$
  \ELSIF {$x_j < x_j^{min}$}
    \STATE $z \gets -2$
  \ENDIF
\ELSIF {$z = 2$}
  \IF {$K_p x_k + x_1 < x_j^{max}$}
    \STATE $z \gets 1$
  \ENDIF
\ELSIF {$z = 1$}
  \IF {$K_p x_k + x_1 > x_j^{max}$}
    \STATE $z \gets 2$
  \ELSIF {$K_p x_k + x_2 < x_j^{max}$ \;\;\;(*)} 
    \STATE $z \gets 0$
  \ENDIF
\ELSIF {$z = -2$}
  \IF {$K_p x_k + x_1 > x_j^{min}$}
    \STATE $z \gets -1$
  \ENDIF
\ELSIF {$z = -1$}
  \IF {$K_p x_k + x_1 < x_j^{min}$}
    \STATE $z \gets -2$
  \ELSIF {$K_p x_k + x_2 > x_j^{min}$ \;\;\;(*)} 
    \STATE $z \gets 0$
  \ENDIF  
\ENDIF
\end{algorithmic}
\end{minipage}

Initialization of the internal state and the discrete variable: 

\hspace*{20ex}\begin{minipage}[t]{15cm}
\begin{algorithmic} 
\IF {$x_j \ge x_j^{max}$}
   \STATE $z \gets 2$
   \STATE $x_1 \gets x_j^{max} - K_p x_k$
   \STATE $x_2 \gets x_1$    
\ELSIF {$x_j \le x_j^{min}$}
   \STATE $z \gets -2$
   \STATE $x_1 \gets x_j^{min} - K_p x_k$
   \STATE $x_2 \gets x_1$   
\ELSE
   \STATE $z \gets 0$
   \STATE $x_1 \gets x_j$
   \STATE $x_2 \gets x_1$  
\ENDIF
\end{algorithmic}
\end{minipage}
   